\subsubsection{Pentadic representation of the microfibre and its incidence}
\label{pi-estrella:desarrollo}

The arrangement of the five regions displays simultaneously
the unity of the ternary quotient and the differentiation retained by the
catalogue. The relevant map is the restriction
\[
 \Psi\big|_{\mathscr R_\pi^{\rm micro}}:
 \mathscr R_\pi^{\rm micro}\longrightarrow\{010211\},
 \qquad \Psi(U)=(-U)\bmod3.
\]
Its five preimages are the \(U_6\) vectors in the preceding table.
Figure \ref{pi-estrella:figura} recovers their radial organization,
preserving the chart order, signatures and multiplicities of the
catalogue (see \ref{pi-estrella:figura}).

% Orden de carta y datos recuperados de fig_cinco_regiones_pi.tex.
\begin{figure}[H]
\centering
\begin{tikzpicture}[x=1cm,y=1cm,font=\small,>=Latex]
 \coordinate (a) at (90:2.2);
 \coordinate (b) at (18:2.2);
 \coordinate (c) at (-54:2.2);
 \coordinate (d) at (-126:2.2);
 \coordinate (e) at (162:2.2);
 \draw[black!35,line width=.45pt]
    (a)--(c)--(e)--(b)--(d)--cycle;
 \node[draw,fill=white,circle,minimum size=1.6cm,align=center,
    inner sep=3pt] (z) at (0,0) {$w_\pi$\\$010211$};
 \foreach \v in {a,b,c,d,e}{
    \fill (\v) circle (2pt);
    \draw[->,line width=.6pt] (\v)--(z);
 }
 \node[above=3mm,align=center] at (a)
   {$U_\pi^\perp$\\$555555$\\multiplicity $432$};
 \node[right=3mm,align=left] at (b)
   {$U_{\pi,1}^{++}$\\$339555$\\multiplicity $144$};
 \node[below right=2mm and 1mm,align=left] at (c)
   {$U_{\pi,2}^{++}$\\$393555$\\multiplicity $144$};
 \node[below left=2mm and 1mm,align=right] at (d)
   {$U_{\pi,3}^{++}$\\$933555$\\multiplicity $144$};
 \node[left=3mm,align=right] at (e)
   {$U_\pi^{--}$\\$933717$\\multiplicity $144$};
 \node[align=center,font=\footnotesize] at (0,-3.55)
   {$1_\perp+3_{++}+1_{--}$\qquad $432+4\cdot144=1008$};
\end{tikzpicture}
\caption{The five regions of \(\pi\) in a pentadic arrangement.
Each position preserves its multiplicative signature and multiplicity;
the radial arrows represent the common projection
\(\Psi(U)=010211\). The star-shaped outline organizes the five
positions of the chart. Their realization in the five Witt octads
is effected through the regional alignment and the
\(4+1\) decomposition in the text.}
\label{pi-estrella:figura}
\end{figure}


Projection to the centre identifies the five ternary images; the
register of outer positions retains the regional information
that makes the incidence realization possible. The distribution
\(1_\perp+3_{++}+1_{--}\) specifies one transverse region, three
positive regions and one negative region. Their normalized weights are
\(3/7\) and \(1/7\) for each of the other four: the
transverse multiplicity is threefold, even though the five images under
\(\Psi\) coincide.

\Needspace{11\baselineskip}
The connection with the Witt star is formulated through the flag
\(B_K\subset H^-_\alpha\), the alignment \(\ell\) and the
lift \(\iota_D\) already constructed. The five regional positions
are realized in the five octads over \(\ell(B_K)\):
\[
\begin{aligned}
 U_\pi^\perp&\longmapsto O^\perp_{\ell(B_K)},\\
 U_\pi^{--}&\longmapsto\iota_D(H^-_\alpha),\\
 \{U_{\pi,1}^{++},U_{\pi,2}^{++},U_{\pi,3}^{++}\}
 &\longmapsto
 \{\iota_D(B_K\cup P):P\in\mathcal P_+\},\\
 \mathcal P_+&=\bigl\{\{1,3\},\{2,11\},\{8,12\}\bigr\}.
\end{aligned}
\]
The three positive positions are assigned to the three pairs
in an ordered chart. The transverse and negative rows
are distinguished by their regional predicates and by the flag.

The two figures are thereby articulated: Figure
\ref{exc:fig-estrella} displays the four hexads of the small Witt
design over \(B_K\); Figure \ref{pi-estrella:figura} displays
the five regions that, under the alignment, occupy the four residual
octads and the transverse octad of the binary design. The
\(4+1\) law, proved in Proposition
\ref{exc:residuo-binario}, specifies the change of incidence. In this
composition, \(K\) determines the marked tetrad and the signature of
\(\alpha\) distinguishes one of its hexads. In addition to their
common ternary value, the regions retain the structure required for this
transport to the exceptional realization.
